% GENERATED FILE. Edit canonical lessons, then run npm run generate:books.
% canonical-source-hash: fnv1a64:6c65ecb2c84918c2
% canonical-lessons: IT-C460-radura, IT-C460-ruscello, IT-C460-torrente, IT-C460-cespuglio, IT-C460-guadare, IT-R460-ripasso-453-460

\chapter{Nel bosco}
\label{ch:it-nel-bosco}
\hlchaptermodality{460}

\begin{chapteropening}
\textbf{By the end of this chapter:} \emph{I can describe a walk in the woods: a clearing, a brook, a mountain stream, a bush and wading across.}

Complete the last lesson of chapter 460: I can describe a walk in the woods: a clearing, a brook, a mountain stream, a bush and wading across.
\end{chapteropening}

\section[la radura]{la radura — the clearing (in a wood)}

\label{lesson:IT-C460-radura}



\begin{quote}
Before the new one: say the Italian for the bank transfer, then the Italian for the instalment.
\end{quote}

\subsection*{You'll want to know: la radura}
\textbf{la radura} — \textquotedblleft{}the clearing (in a wood)\textquotedblright{}.

\begin{quote}
Abbiamo mangiato in una radura.
\end{quote}

\begin{quote}
\emph{We ate in a clearing.}
\end{quote}

Where the trees are \emph{radi}, far apart.

\begin{cousinweb}
\emph{radura} is built on \emph{rado} 'sparse', which is inherited from Latin \emph{rārus} 'thin, far apart' (EtymDB).
\end{cousinweb}

\begin{grammarlens}[title={What you've built}]
One more word for telling, explaining and describing at B1.
\end{grammarlens}

\subsection*{Your turn}
\begin{itemize}
\raggedright
  \item \emph{Say it:} \emph{la radura}
  \item \emph{Say it:} \emph{la radura}, once more
  \item \emph{Say it:} a sentence of your own with \emph{la radura}
  \item \emph{Recall:} say the Italian for the bank transfer, then the Italian for the instalment, then say \emph{la radura} again
\end{itemize}

\begin{tcolorbox}[breakable,colback=teal!4,colframe=teal!35!black,title={Before you move on}]
What does la radura mean? (\textquotedblleft{}The clearing (in a wood)\textquotedblright{}.) Say it once more.
\end{tcolorbox}
\section[il ruscello]{il ruscello — the brook, the stream}

\label{lesson:IT-C460-ruscello}



\begin{quote}
Before the new one: say the Italian for the instalment, then the Italian for the clearing.
\end{quote}

\subsection*{You'll want to know: il ruscello}
\textbf{il ruscello} — \textquotedblleft{}the brook, the stream\textquotedblright{}.

\begin{quote}
Beviamo l'acqua fresca del ruscello.
\end{quote}

\begin{quote}
\emph{We drink the cool water of the brook.}
\end{quote}

A small, quiet stream; \emph{il fiume} is a river.

\begin{cousinweb}
\emph{ruscello} is inherited from Vulgar Latin \emph{rīvuscellus} (EtymDB), a little form of Latin \emph{rīvus} 'stream'.
\end{cousinweb}

\begin{grammarlens}[title={What you've built}]
One more word for telling, explaining and describing at B1.
\end{grammarlens}

\subsection*{Your turn}
\begin{itemize}
\raggedright
  \item \emph{Say it:} \emph{il ruscello}
  \item \emph{Say it:} \emph{il ruscello}, once more
  \item \emph{Say it:} a sentence of your own with \emph{il ruscello}
  \item \emph{Recall:} say the Italian for the instalment, then the Italian for the clearing, then say \emph{il ruscello} again
\end{itemize}

\begin{tcolorbox}[breakable,colback=teal!4,colframe=teal!35!black,title={Before you move on}]
What does il ruscello mean? (\textquotedblleft{}The brook, the stream\textquotedblright{}.) Say it once more.
\end{tcolorbox}
\section[il torrente]{il torrente — the mountain stream, the torrent}

\label{lesson:IT-C460-torrente}



\begin{quote}
Before the new one: say the Italian for the clearing, then the Italian for the brook.
\end{quote}

\subsection*{You'll want to know: il torrente}
\textbf{il torrente} — \textquotedblleft{}the mountain stream, the torrent\textquotedblright{}.

\begin{quote}
Dopo la pioggia il torrente è pieno.
\end{quote}

\begin{quote}
\emph{After the rain the mountain stream is full.}
\end{quote}

Fast, and often dry in summer.

\begin{cousinweb}
\emph{torrente} is inherited from Latin \emph{torrēns, torrentem} 'rushing stream' (EtymDB), a form of \emph{torrēre} 'to scorch'.
\end{cousinweb}

\begin{grammarlens}[title={What you've built}]
One more word for telling, explaining and describing at B1.
\end{grammarlens}

\subsection*{Your turn}
\begin{itemize}
\raggedright
  \item \emph{Say it:} \emph{il torrente}
  \item \emph{Say it:} \emph{il torrente}, once more
  \item \emph{Say it:} a sentence of your own with \emph{il torrente}
  \item \emph{Recall:} say the Italian for the clearing, then the Italian for the brook, then say \emph{il torrente} again
\end{itemize}

\begin{tcolorbox}[breakable,colback=teal!4,colframe=teal!35!black,title={Before you move on}]
What does il torrente mean? (\textquotedblleft{}The mountain stream, the torrent\textquotedblright{}.) Say it once more.
\end{tcolorbox}
\section[il cespuglio]{il cespuglio — the bush, the shrub}

\label{lesson:IT-C460-cespuglio}



\begin{quote}
Before the new one: say the Italian for the brook, then the Italian for the mountain stream.
\end{quote}

\subsection*{You'll want to know: il cespuglio}
\textbf{il cespuglio} — \textquotedblleft{}the bush, the shrub\textquotedblright{}.

\begin{quote}
Il gatto si nasconde dietro un cespuglio.
\end{quote}

\begin{quote}
\emph{The cat hides behind a bush.}
\end{quote}

\emph{Un cespuglio di rose} — a rose bush.

\begin{cousinweb}
\emph{cespuglio} is traced in EtymDB to Latin \emph{caespuculum}, a little form of \emph{caespes} 'turf, sod'; it is inherited.
\end{cousinweb}

\begin{grammarlens}[title={What you've built}]
One more word for telling, explaining and describing at B1.
\end{grammarlens}

\subsection*{Your turn}
\begin{itemize}
\raggedright
  \item \emph{Say it:} \emph{il cespuglio}
  \item \emph{Say it:} \emph{il cespuglio}, once more
  \item \emph{Say it:} a sentence of your own with \emph{il cespuglio}
  \item \emph{Recall:} say the Italian for the brook, then the Italian for the mountain stream, then say \emph{il cespuglio} again
\end{itemize}

\begin{tcolorbox}[breakable,colback=teal!4,colframe=teal!35!black,title={Before you move on}]
What does il cespuglio mean? (\textquotedblleft{}The bush, the shrub\textquotedblright{}.) Say it once more.
\end{tcolorbox}
\section[guadare]{guadare — to wade across, to ford}

\label{lesson:IT-C460-guadare}



\begin{quote}
Before the new one: say the Italian for the mountain stream, then the Italian for the bush.
\end{quote}

\subsection*{You'll want to know: guadare}
\textbf{guadare} — \textquotedblleft{}to wade across, to ford\textquotedblright{}.

\begin{quote}
Abbiamo guadato il ruscello a piedi.
\end{quote}

\begin{quote}
\emph{We waded across the brook on foot.}
\end{quote}

\emph{Il guado} — the ford.

\begin{cousinweb}
\emph{guadare} is built on \emph{guado} 'ford', which is inherited from Latin \emph{vadum} 'shallow place' (Lewis \& Short), its \emph{gu-} shaped by a Germanic word.
\end{cousinweb}

\begin{grammarlens}[title={What you've built}]
One more word for telling, explaining and describing at B1.
\end{grammarlens}

\subsection*{Your turn}
\begin{itemize}
\raggedright
  \item \emph{Say it:} \emph{guadare}
  \item \emph{Say it:} \emph{guadare}, once more
  \item \emph{Say it:} a sentence of your own with \emph{guadare}
  \item \emph{Recall:} say the Italian for the mountain stream, then the Italian for the bush, then say \emph{guadare} again
\end{itemize}

\begin{tcolorbox}[breakable,colback=teal!4,colframe=teal!35!black,title={Before you move on}]
What does guadare mean? (\textquotedblleft{}To wade across, to ford\textquotedblright{}.) Say it once more.
\end{tcolorbox}
\section[(dialogue)]{Ripasso dei capitoli 453–460}

\label{lesson:IT-R460-ripasso-453-460}



\begin{quote}
Say the Italian for the hardware shop and the haberdashery.
\end{quote}

\subsection*{Your turn}
Say these aloud:

\begin{itemize}
\raggedright
  \item the hardware shop, the haberdashery, the perfume shop, the herbalist's, the jeweller's
  \item the delicatessen, the fishmonger's, the greengrocer, the cake shop, the rotisserie
  \item the cobbler, to resole, the tailor, to shorten, the locksmith
  \item the hairdresser, the barber, the beautician, the optician, to remove stains from
  \item the cashier, the small change, the fitting room, the clearance sale, to haggle
  \item the registered letter, the shipping, the sender, to put postage on, to deliver
  \item to withdraw, the withdrawal, to charge, the bank transfer, the instalment
  \item the clearing, the brook, the mountain stream, the bush, to wade across
\end{itemize}

\begin{tcolorbox}[breakable,colback=teal!4,colframe=teal!35!black,title={Before you move on}]
The hardware shop? (\textbf{la ferramenta}.) The haberdashery? (\textbf{la merceria}.) The perfume shop? (\textbf{la profumeria}.) The herbalist's? (\textbf{l'erboristeria}.) The jeweller's? (\textbf{la gioielleria}.) The delicatessen? (\textbf{la salumeria}.) The fishmonger's? (\textbf{la pescheria}.) The greengrocer? (\textbf{il fruttivendolo}.) The cake shop? (\textbf{la pasticceria}.) The rotisserie? (\textbf{la rosticceria}.) The cobbler? (\textbf{il calzolaio}.) To resole? (\textbf{risuolare}.) The tailor? (\textbf{il sarto}.) To shorten? (\textbf{accorciare}.) The locksmith? (\textbf{il fabbro}.) The hairdresser? (\textbf{il parrucchiere}.) The barber? (\textbf{il barbiere}.) The beautician? (\textbf{l'estetista}.) The optician? (\textbf{l'ottico}.) To remove stains from? (\textbf{smacchiare}.) The cashier? (\textbf{il cassiere}.) The small change? (\textbf{gli spiccioli}.) The fitting room? (\textbf{il camerino}.) The clearance sale? (\textbf{la svendita}.) To haggle? (\textbf{mercanteggiare}.) The registered letter? (\textbf{la raccomandata}.) The shipping? (\textbf{la spedizione}.) The sender? (\textbf{il mittente}.) To put postage on? (\textbf{affrancare}.) To deliver? (\textbf{recapitare}.) To withdraw? (\textbf{prelevare}.) The withdrawal? (\textbf{il prelievo}.) To charge? (\textbf{addebitare}.) The bank transfer? (\textbf{il bonifico}.) The instalment? (\textbf{la rata}.) The clearing? (\textbf{la radura}.) The brook? (\textbf{il ruscello}.) The mountain stream? (\textbf{il torrente}.) The bush? (\textbf{il cespuglio}.) To wade across? (\textbf{guadare}.)
\end{tcolorbox}
